\documentclass[12pt]{article}

\usepackage{amsmath}
\usepackage{amssymb}
\usepackage{geometry}
\usepackage{setspace}
\usepackage{booktabs}
\usepackage{algorithm}
\usepackage{algpseudocode}

\geometry{margin=1in}
\onehalfspacing

\title{Advanced Optimization Formulation \\ Examination Room Scheduling}
\author{}
\date{}

\begin{document}

\maketitle

\section{Notation}

\subsection*{Sets}

\begin{itemize}
\item $P$ : set of healthcare providers
\item $D$ : set of days in the planning horizon
\item $A$ : set of appointments
\item $A(p,d)$ : appointments of provider $p$ on day $d$, sorted by start time
\item $R$ : set of examination rooms (ER1--ER16)
\item $C_p \subseteq R$ : cluster of rooms available to provider $p$ (rooms within proximity threshold $\delta_{\max}$)
\item $S(p,d)$ : set of feasible schedules for provider $p$ on day $d$
\item $O$ : set of appointment pairs $(a,b)$ whose time intervals overlap (with robust buffer)
\item $K(p,d)$ : set of consecutive appointment pairs for provider $p$ on day $d$, i.e., $(a_i, a_{i+1})$ ordered by start time
\end{itemize}

\subsection*{Parameters}

\begin{itemize}
\item $t_a$ : start time of appointment $a$ (minutes from midnight, fixed)
\item $d_a$ : duration of appointment $a$ (minutes)
\item $\Delta_a$ : robust duration buffer for appointment $a$; $\Delta_a = \delta_{\mathrm{frac}} \cdot d_a$
\item $dist_{rr'}$ : travel distance (metres) between rooms $r$ and $r'$
\item $\gamma$ : penalty per room switch ($\gamma = 10$)
\item $\eta$ : penalty weight per metre of travel ($\eta = 1$)
\item $\rho_a$ : no-show probability for the provider of appointment $a$
\item $M$ : large penalty for room capacity violation ($M = 10^{4}$)
\item $c_s$ : total cost of schedule $s$ (sum of switch and travel penalties)
\item $\alpha_{as} \in \{0,1\}$ : 1 if appointment $a$ is covered by schedule $s$
\item $\beta_{rts} \in \{0,1\}$ : 1 if room $r$ is occupied at minute $t$ under schedule $s$
\end{itemize}

\subsection*{Overlap Definition}

Two appointments $(a,b)$ are in $O$ (i.e., conflict) if their buffered intervals overlap:

\[
t_a < t_b + d_b + \Delta_b \quad \text{and} \quad t_b < t_a + d_a + \Delta_a
\]

\newpage

\section{Model 1: Feasibility Packing Model}

Model~1 assigns all appointments to examination rooms while respecting time conflicts.
The objective is to minimise the number of distinct rooms opened, yielding compact schedules.

\subsection*{Decision Variables}

\[
x_{ar} =
\begin{cases}
1 & \text{if appointment } a \text{ is assigned to room } r \\
0 & \text{otherwise}
\end{cases}
\qquad \forall a \in A,\; r \in R
\]

\[
y_r =
\begin{cases}
1 & \text{if room } r \text{ is used by at least one appointment} \\
0 & \text{otherwise}
\end{cases}
\qquad \forall r \in R
\]

\subsection*{Objective Function}

\[
\min \sum_{r \in R} y_r
\]

Minimising the number of rooms opened encourages compact, clinically practical assignments.
(Note: the alternative objective $\min \sum_{a}\sum_{r} x_{ar}$ is trivially constant at $|A|$ due
to the assignment equality constraint, so $y_r$ is used instead.)

\subsection*{Constraints}

\subsubsection*{Appointment Assignment}

\[
\sum_{r \in R} x_{ar} = 1 \quad \forall a \in A
\]

Each appointment is assigned to exactly one room.

\subsubsection*{Room Conflict Constraint}

\[
x_{ar} + x_{br} \le 1 \quad \forall (a,b) \in O,\; \forall r \in R
\]

Appointments with overlapping (buffered) intervals cannot share a room.

\subsubsection*{Room Open Linking Constraint}

\[
x_{ar} \le y_r \quad \forall a \in A,\; \forall r \in R
\]

A room is marked open whenever any appointment is assigned to it.

\subsubsection*{Binary Domain}

\[
x_{ar} \in \{0,1\}, \quad y_r \in \{0,1\}
\]

\newpage

\section{Model 2: Provider-Day Schedule Generation}

For each provider $p$ and day $d$, Model~2 assigns appointments to rooms while minimising
room switching and provider travel. Because appointment start times $t_a$ are fixed by the
booking system, time is not a decision variable; the index $t$ is dropped and
$x_{ar}$ fully determines the assignment.

\subsection*{Decision Variables}

\[
x_{ar} =
\begin{cases}
1 & \text{if appointment } a \in A(p,d) \text{ is assigned to room } r \in C_p \\
0 & \text{otherwise}
\end{cases}
\]

To track room transitions between consecutive appointments, define the linearisation variable:

\[
w_{ab r r'} =
\begin{cases}
1 & \text{if consecutive pair } (a,b) \in K(p,d) \text{ uses rooms } r \text{ then } r' \\
0 & \text{otherwise}
\end{cases}
\quad \forall (a,b) \in K(p,d),\; r,r' \in C_p
\]

\subsection*{Objective Function}

\[
\min \;
\gamma \sum_{(a,b)\in K} \sum_{r \in C_p} \sum_{\substack{r' \in C_p \\ r' \neq r}} w_{ab r r'}
\;+\;
\eta \sum_{(a,b)\in K} \sum_{r \in C_p} \sum_{\substack{r' \in C_p \\ r' \neq r}} dist_{rr'}\, w_{ab r r'}
\]

The objective penalises room switches (count) and total provider travel distance.

\subsection*{Constraints}

\subsubsection*{Appointment Assignment}

\[
\sum_{r \in C_p} x_{ar} = 1 \quad \forall a \in A(p,d)
\]

\subsubsection*{Room Conflict Constraint}

\[
x_{ar} + x_{br} \le 1 \quad \forall (a,b) \in O,\; \forall r \in C_p
\]

\subsubsection*{Linearisation of Transition Variables}

For each consecutive pair $(a,b) \in K(p,d)$ and all $r, r' \in C_p$:

\[
w_{ab r r'} \le x_{ar} \tag{L1}
\]
\[
w_{ab r r'} \le x_{br'} \tag{L2}
\]
\[
w_{ab r r'} \ge x_{ar} + x_{br'} - 1 \tag{L3}
\]

These three constraints enforce $w_{ab r r'} = x_{ar} \cdot x_{br'}$ without
introducing nonlinearity.

\subsubsection*{Exactly One Transition per Consecutive Pair}

\[
\sum_{r \in C_p} \sum_{r' \in C_p} w_{ab r r'} = 1 \quad \forall (a,b) \in K(p,d)
\]

\subsubsection*{Room Cluster Constraint}

\[
x_{ar} = 0 \quad \text{if } r \notin C_p
\]

Providers are restricted to rooms within proximity threshold $\delta_{\max}$ of their
assigned home room.

\subsubsection*{Binary Domain}

\[
x_{ar} \in \{0,1\}, \quad w_{ab r r'} \in \{0,1\}
\]

\newpage

\section{Model 3: Master Schedule Selection}

Model~3 selects one schedule per provider-day from the pool generated by Model~2
such that every appointment is covered and no room is double-booked.
Room capacity is enforced as a soft constraint with a large penalty $M$ to guarantee
feasibility when residual cross-provider conflicts exist.

\subsection*{Decision Variables}

\[
\lambda_s =
\begin{cases}
1 & \text{if schedule } s \in S(p,d) \text{ is selected} \\
0 & \text{otherwise}
\end{cases}
\]

\[
\sigma_{rt} \ge 0 \quad \text{(slack for room } r \text{ at minute } t\text{)}
\]

\subsection*{Objective Function}

\[
\min \sum_{p \in P} \sum_{d \in D} \sum_{s \in S(p,d)} c_s \lambda_s
\;+\;
M \sum_{(r,t)} \sigma_{rt}
\]

where $M = 10^{4}$ heavily penalises any cross-provider room conflict.

\subsection*{Constraints}

\subsubsection*{Appointment Coverage}

\[
\sum_{p \in P}\sum_{d \in D} \sum_{s \in S(p,d)} \alpha_{as}\,\lambda_s = 1
\quad \forall a \in A
\]

Each appointment is covered by exactly one selected schedule.

\subsubsection*{Room Capacity (Soft)}

\[
\sum_{p \in P} \sum_{s \in S(p,d)} \beta_{rts}\,\lambda_s - \sigma_{rt} \le 1
\quad \forall r \in R,\; t \in T,\; d \in D
\]

At most one provider may occupy room $r$ at time $t$; the slack $\sigma_{rt}$
absorbs residual conflicts at cost $M$ per unit.

\subsubsection*{One Schedule per Provider-Day}

\[
\sum_{s \in S(p,d)} \lambda_s = 1 \quad \forall p \in P,\; \forall d \in D
\]

\subsubsection*{Binary and Non-negativity Domain}

\[
\lambda_s \in \{0,1\}, \quad \sigma_{rt} \ge 0
\]

\newpage

\section{Robust Scheduling Extension}

\subsection*{Duration Buffer}

To protect against appointment overruns, a robust buffer
$\Delta_a = \delta_{\mathrm{frac}} \cdot d_a$ is added to each appointment when
computing overlaps in Models~1 and~2.
Two appointments $(a,b)$ are treated as conflicting whenever:

\[
t_a < t_b + d_b + \Delta_b \quad \text{and} \quad t_b < t_a + d_a + \Delta_a
\]

In the implemented policies, $\delta_{\mathrm{frac}} = 0.10$ (10\% buffer).

\subsection*{No-Show Overbooking Adjustment}

For Policy~D, effective appointment durations are reduced by the provider's
historical no-show rate before conflict detection, allowing the schedule to
accommodate additional appointments:

\[
\hat{d}_a = (1 - \rho_a)\,d_a
\]

where $\rho_a$ is estimated from historical data (minimum 5 observations;
otherwise the clinic-wide average is used). The reduced duration $\hat{d}_a$
replaces $d_a$ in the overlap computation for Models~1 and~2, increasing
effective room capacity.

\newpage

\section{Implemented Scheduling Policies}

Four policies are evaluated using the three-model framework:

\begin{center}
\begin{tabular}{llcc}
\toprule
\textbf{Policy} & \textbf{Description} & \textbf{Proximity $\delta_{\max}$ (m)} & \textbf{Buffer $\delta_{\mathrm{frac}}$} \\
\midrule
A & Single room per provider (baseline) & 0 & 0 \\
B & Cluster of nearby rooms & 4 & 0 \\
C & Cluster + robust duration buffer & 4 & 0.10 \\
D & Cluster + buffer + no-show overbooking & 4 & 0.10 \\
\bottomrule
\end{tabular}
\end{center}

\subsection*{Policy A (Baseline)}

Each provider is assigned their pre-designated single room for the entire day.
No optimisation is performed; this mirrors current clinical practice and serves
as the coverage and utilisation benchmark.

\subsection*{Policies B--D (Optimised)}

Model~2 is solved sequentially per provider (busiest provider first) with
Model~2-with-exclusions ensuring no two providers claim the same (room, time) slot.
Model~3 then selects the globally optimal combination of provider-day schedules.

\section{Column Generation Algorithm}

The three models are integrated within a \textbf{Dantzig-Wolfe column generation}
framework. Model~2 is the \emph{pricing subproblem} and Model~3 is the
\emph{Restricted Master Problem} (RMP). The algorithm iterates between them until
no improving column exists.

\subsection*{Reduced Cost}

Given dual variables $\pi_a \geq 0$ (appointment coverage) and $\mu_{rt} \geq 0$
(room capacity) from the LP relaxation of Model~3, the \textbf{reduced cost} of a
candidate schedule $s$ is:

\[
\bar{c}_s = c_s - \sum_{a \in A} \pi_a \,\alpha_{as}
            - \sum_{(r,t)} \mu_{rt} \,\beta_{rts}
\]

A schedule with $\bar{c}_s < 0$ improves the master problem objective and should be
added as a new column.

\subsection*{Algorithm}

\begin{algorithm}[H]
\caption{Column Generation for Examination Room Scheduling}
\begin{algorithmic}[1]
\State \textbf{Initialise:} Run Model~2 once per (provider, day) to build initial
       column pool $\mathcal{S}$.
\Repeat
    \State Solve \textbf{LP relaxation} of Model~3 over current $\mathcal{S}$.
    \State Extract dual variables $\pi_a$ (appointment) and $\mu_{rt}$ (room capacity).
    \For{each (provider $p$, day $d$)}
        \State Solve Model~2 with modified objective using reduced costs:
        \[
          \min \; c_s - \sum_{a} \pi_a \,\alpha_{as} - \sum_{(r,t)} \mu_{rt}\,\beta_{rts}
        \]
        \If{$\bar{c}_s < 0$}
            \State Add new schedule $s$ to column pool $\mathcal{S}$.
        \EndIf
    \EndFor
\Until{no new column has $\bar{c}_s < 0$} \Comment{LP optimality condition}
\State Solve \textbf{ILP} of Model~3 over final $\mathcal{S}$ to obtain integer solution.
\end{algorithmic}
\end{algorithm}

\noindent The LP relaxation at termination provides a \textbf{lower bound}; the ILP
solution provides an upper bound. Their gap measures solution quality.

\newpage

\section{Lagrangian Relaxation}

Lagrangian Relaxation (LR) is applied to Model~3 to obtain a tractable lower bound
by dualising the \textbf{appointment coverage constraints}.

\subsection*{Lagrangian Relaxation of Model 3}

Let $u_a \geq 0$ be the Lagrange multiplier for appointment $a$'s coverage
constraint. The Lagrangian relaxation of Model~3 is:

\[
L(\mathbf{u}) = \min_{\lambda} \left[
    \sum_{s} c_s \lambda_s
    + M \sum_{(r,t)} \sigma_{rt}
    + \sum_{a \in A} u_a \!\left(1 - \sum_{s} \alpha_{as} \lambda_s \right)
\right]
\]

subject to the room capacity and one-schedule-per-provider-day constraints, and
$\lambda_s \in \{0,1\}$. Expanding:

\[
L(\mathbf{u}) = \sum_{a} u_a
+ \min_{\lambda} \left[
    \sum_{s} \left( c_s - \sum_{a} u_a \alpha_{as} \right) \lambda_s
    + M \sum_{(r,t)} \sigma_{rt}
\right]
\]

subject to room capacity and one-schedule-per-provider-day constraints.

\subsection*{Lagrangian Dual}

$L(\mathbf{u})$ provides a lower bound on the optimal objective for any $\mathbf{u} \geq 0$.
The \textbf{Lagrangian dual} seeks the tightest such bound:

\[
z_{LD} = \max_{\mathbf{u} \geq 0} \; L(\mathbf{u})
\]

\subsection*{Subgradient Method}

The Lagrangian dual is solved iteratively using the \textbf{subgradient method}.
At iteration $k$:

\begin{enumerate}
    \item Solve the Lagrangian subproblem to obtain $\lambda^k$.
    \item Compute the subgradient for each appointment $a$:
    \[
        g_a^k = 1 - \sum_{s} \alpha_{as} \lambda_s^k
    \]
    \item Update multipliers:
    \[
        u_a^{k+1} = \max\!\left(0,\; u_a^k - t_k \cdot g_a^k \right)
    \]
    \item Update step size using the Polyak formula:
    \[
        t_k = \frac{\lambda^k \left( z_{UB} - L(\mathbf{u}^k) \right)}{\| \mathbf{g}^k \|^2}
    \]
    where $z_{UB}$ is the best known feasible (upper bound) objective value and
    $\lambda^k \in (0, 2]$ is a scalar (typically initialised at $2$ and halved
    when $L(\mathbf{u})$ fails to improve for several consecutive iterations).
\end{enumerate}

\noindent The algorithm terminates when $\|\mathbf{g}^k\| \approx 0$ (coverage
constraints satisfied), the duality gap $z_{UB} - z_{LD}$ is sufficiently small,
or the iteration limit is reached.

\newpage

\section{Tabu Search Metaheuristic}

Following the exact optimisation phase, a \textbf{Tabu Search} (TS) metaheuristic
is applied to improve the schedule by escaping local optima that the MIP solver
may not reach within its time limit.

\subsection*{Solution Representation}

A solution is a complete room assignment:
$\mathbf{x} = \{ (a, r_a) : a \in A \}$, where $r_a \in C_p$ is the room
assigned to appointment $a$ of provider $p$.

\subsection*{Neighbourhood Structure}

The neighbourhood $N(\mathbf{x})$ is defined by \textbf{single-appointment room
swaps}: for each appointment $a$, replace $r_a$ with any other room $r' \in C_p$
such that $r' \neq r_a$ and assigning $r'$ to $a$ does not create a conflict
with another appointment already in $r'$.

\subsection*{Objective and Penalty}

The TS objective augments the scheduling cost with a penalty for constraint violations:

\[
f(\mathbf{x}) = \gamma \cdot \text{switches}(\mathbf{x})
              + \eta \sum_{(a,b)\text{ consecutive}} dist_{r_a r_b}
              + P \cdot \text{conflicts}(\mathbf{x})
\]

where $P \gg \gamma$ is a large violation penalty, and $\text{conflicts}(\mathbf{x})$
counts the number of (room, time) double-bookings.

\subsection*{Algorithm}

\begin{algorithm}[H]
\caption{Tabu Search for Room Assignment Improvement}
\begin{algorithmic}[1]
\State \textbf{Initialise:} $\mathbf{x}_0 \leftarrow$ ILP solution from Model~3;\;
       $\mathbf{x}^* \leftarrow \mathbf{x}_0$;\; Tabu list $\Gamma \leftarrow \emptyset$
\For{$k = 1, 2, \ldots, K_{\max}$}
    \State Generate candidate set $\mathcal{C} \subseteq N(\mathbf{x}_{k-1})$
    \State Select best $\mathbf{x}' \in \mathcal{C}$ such that
           $(a, r') \notin \Gamma$ \textbf{or} $f(\mathbf{x}') < f(\mathbf{x}^*)$
           \Comment{Aspiration criterion}
    \State $\mathbf{x}_k \leftarrow \mathbf{x}'$
    \State Add move $(a, r_a \to r')$ to $\Gamma$; if $|\Gamma| > k_{\max}$ remove oldest entry
    \If{$f(\mathbf{x}_k) < f(\mathbf{x}^*)$}
        \State $\mathbf{x}^* \leftarrow \mathbf{x}_k$ \Comment{Update best solution}
    \EndIf
\EndFor
\State \Return $\mathbf{x}^*$
\end{algorithmic}
\end{algorithm}

\noindent \textbf{Parameters:} tabu list size $k_{\max} \in [5, 20]$;
iteration limit $K_{\max}$; aspiration criterion — override tabu status if the
move produces a solution better than the current global best $\mathbf{x}^*$.

\section{Integrated Solution Framework}

The complete solution procedure proceeds as follows:

\begin{enumerate}
\item \textbf{Model~1} (Feasibility Packing, solved as ILP via branch-and-bound):
      verify all appointments can be packed into available rooms; provides an initial
      compact room assignment and confirms problem feasibility.

\item \textbf{Column Generation} (Models~2 and~3 iterated):
      Model~2 generates provider-day schedules as pricing subproblems;
      Model~3 (Restricted Master Problem) selects the globally optimal combination.
      The LP relaxation of Model~3 yields a \textbf{lower bound}; the ILP gives
      the integer solution.

\item \textbf{Lagrangian Relaxation} (applied to Model~3):
      dualise appointment coverage constraints; solve the Lagrangian subproblem
      iteratively via the subgradient method to obtain a complementary lower bound
      and assess the duality gap.

\item \textbf{Tabu Search} (post-optimisation improvement):
      starting from the column generation integer solution, apply TS to improve
      room assignments by exploring the single-swap neighbourhood while maintaining
      a tabu list to prevent cycling.
\end{enumerate}

\end{document}
